% Einheit 2 von 3 – Kathete und Umkehrung (bank/pythagoras/e2.jsonl, zusammenbau v0.1)
%% TODO Zweigzeile Teil 1: was hier gelernt wird (Satz in Schülersprache aus dem Einheitstitel) – Einheit 2
\einheitenkopf[e2]{Einheit 2 von 3 · Kathete und Umkehrung}
\zweigzeile{ab Kl. 8, je nach Buch bis Kl. 9 · P10 oft}
\setcounter{aufgabe}{16}

%% TODO Titel: Ich-kann-Satz fehlt in der Bank (Platzhalter: Kettenname) – Nr. 17
%% TODO Anweisung: ein Satz über der ersten Teilaufgabe fehlt in der Bank – Nr. 17
\begin{aufgabe}{Welche Seite ist die längste?}
\begin{teile}
\teil Welche Seite kann die Hypotenuse sein? \\ \kreuz{$4$ cm} \\ \kreuz{$8{,}5$ cm} \\ \kreuz{$7{,}5$ cm}
\end{teile}
\end{aufgabe}

%% TODO Titel: Ich-kann-Satz fehlt in der Bank (Platzhalter: Kettenname) – Nr. 18
%% TODO Anweisung: ein Satz über der ersten Teilaufgabe fehlt in der Bank – Nr. 18
\begin{aufgabe}{Kathete}
\begin{teile}
\teil Leiterlänge und Abstand des Leiterfußes zur Wand sind gegeben, gesucht ist die Höhe an der Wand. Plus oder minus? \\ \kreuz{Hypotenuse gesucht – plus} \\ \kreuz{Kathete gesucht – minus}
\end{teile}
%% TODO Darstellung für schwach fehlt in der Bank (pythagoras-e2-k2-s1-v1; Abschnitt „Für schwache Schüler“)
\swz{Hypotenuse $c = 89$ cm, Kathete $b = 80$ cm – Kathete $a$?}{}
%% TODO Darstellung für schwach fehlt in der Bank (pythagoras-e2-k2-s1-v2; Abschnitt „Für schwache Schüler“)
\swz{Hypotenuse $c = 97$ m, Kathete $a = 65$ m – Kathete $b$?}{}
%% TODO Darstellung für schwach fehlt in der Bank (pythagoras-e2-k2-s1-v3; Abschnitt „Für schwache Schüler“)
\swz{Hypotenuse $c = 73$ mm, Kathete $b = 48$ mm – Kathete $a$?}{}
%% TODO Darstellung für schwach fehlt in der Bank (pythagoras-e2-k2-s1-v4; Abschnitt „Für schwache Schüler“)
\swz{Hypotenuse $c = 87$ cm, Kathete $a = 60$ cm – Kathete $b$?}{}
%% TODO Darstellung für schwach fehlt in der Bank (pythagoras-e2-k2-s1-v5; Abschnitt „Für schwache Schüler“)
\swz{Hypotenuse $c = 50$ m, Kathete $b = 14$ m – Kathete $a$?}{}
\swfrage{Was bleibt gleich, was ändert sich?}
%% TODO Darstellung für schwach fehlt in der Bank (pythagoras-e2-k2-s2-v1; Abschnitt „Für schwache Schüler“)
\swz{Hypotenuse $c = 11$ cm, Kathete $b = 6$ cm – Kathete $a$ auf eine Dezimale?}{}
%% TODO Darstellung für schwach fehlt in der Bank (pythagoras-e2-k2-s3-v1; Abschnitt „Für schwache Schüler“)
\swz{Hypotenuse $c = 6{,}5$ m, Kathete $b = 2{,}5$ m – Kathete $a$?}{}
%% TODO Darstellung für schwach fehlt in der Bank (pythagoras-e2-k2-s4-v1; Abschnitt „Für schwache Schüler“)
\swz{Hypotenuse $c = 7{,}5$ cm, Kathete $b = 4{,}5$ cm – Kathete $a$? Ist $a$ kürzer als $c$?}{}
\begin{teile}
\teil Welche Gleichung gilt für die Kathete $f$? \\ \kreuz{$f = \sqrt{g^2 - e^2}$} \\ \kreuz{$f = \sqrt{g^2 + e^2}$} \\ \kreuz{$f = g^2 - e^2$} \\ \kreuz{$f = g - e$}
\end{teile}
%% TODO Darstellung für schwach fehlt in der Bank (pythagoras-e2-k2-s6-v1; Abschnitt „Für schwache Schüler“)
\swz[3]{Hypotenuse $8$ cm, Kathete $4{,}9$ cm. Weise nach, dass die andere Kathete etwa $6{,}3$ cm lang ist.}{}
\swz{Ein Waldweg führt $250$ m bergauf und kommt dabei waagerecht $240$ m voran. Wie groß ist der Höhenunterschied?}{\dreieckrw{5}{1.46}{$240$ m}{?}{$250$ m}}
%% TODO Darstellung für schwach fehlt in der Bank (pythagoras-e2-k2-s8-v1; Abschnitt „Für schwache Schüler“)
\swz{Seiten $33$ cm, $56$ cm, $65$ cm – rechtwinklig?}{}
%% TODO Darstellung für schwach fehlt in der Bank (pythagoras-e2-k2-s9-v1; Abschnitt „Für schwache Schüler“)
\swz[3]{Ein dreieckiges Beet hat die Seiten $1{,}8$ m, $2{,}4$ m und $3{,}0$ m. Begründe, dass es einen rechten Winkel hat, und gib an, wo er liegt.}{}
\swz[3]{Eine Seilrutsche ist $426$ m lang und überbrückt waagerecht $390$ m. Berechne den Höhenunterschied zwischen Start und Ziel auf eine Dezimale \hfill (P10 2024 OS).}{\dreieckrw{5}{2.2}{$390$ m}{$h$}{$426$ m}}
\swz[3]{Im Dreieck PQR ist der Winkel bei R stumpf. Die Höhe von R auf die Seite PQ ist $5{,}3$ m lang, ihr Fußpunkt F liegt auf PQ. Es gilt $PR = 9{,}6$ m. Weise nach, dass $PF \approx 8{,}0$ m ist \hfill (P10 2022 OS).}{}
\swz[3]{Ein Fenster hat die Form eines Drachenvierecks ABCD. Seine Diagonalen schneiden sich in M rechtwinklig, und es gilt $AM = MC = DM = 32$ cm. Begründe mit der Umkehrung des Satzes des Pythagoras, dass das Fenster bei D einen rechten Winkel hat \hfill (P10 2025 OS).}{}
\end{aufgabe}

%% TODO Titel: Ich-kann-Satz fehlt in der Bank (Platzhalter: Kettenname) – Nr. 19
%% TODO Anweisung: ein Satz über der ersten Teilaufgabe fehlt in der Bank – Nr. 19
\begin{aufgabe}{pythagoreische Tripel erkennen und vervielfachen}
%% TODO Darstellung für schwach fehlt in der Bank (pythagoras-e2-k3-s1-v1; Abschnitt „Für schwache Schüler“)
\swz{Pythagoreisches Tripel: drei ganze Zahlen, bei denen die Quadrate der zwei kleineren zusammen das Quadrat der größten ergeben. Sind $11$, $60$, $61$ ein Tripel?}{}
\end{aufgabe}

%% TODO Titel: Ich-kann-Satz fehlt in der Bank (Platzhalter: Kettenname) – Nr. 20
%% TODO Anweisung: ein Satz über der ersten Teilaufgabe fehlt in der Bank – Nr. 20
\begin{aufgabe}{Kathete – Fehler finden, Begründen, Anwendung, Darstellungswechsel}
\swz[3]{Hypotenuse $58$ cm, Kathete $40$ cm. Jan rechnet: \rechnung{a^2 &= 58^2 + 40^2 \\ a^2 &= 4\,964 \\ a &\approx 70{,}5 \text{ cm}} Wo ist der Fehler?}{}
\swfrage{Warum wird subtrahiert, wenn eine Kathete gesucht ist?}
\swz[3]{Eine Drehleiter der Feuerwehr ist $30$ m lang ausgefahren. Ihr unteres Ende sitzt $3$ m über dem Boden, $14$ m waagerecht vom Haus entfernt. Reicht sie bis zu einem Fenster in $28$ m Höhe?}{}
\swa{Zeichne ein rechtwinkliges Dreieck, zu dem $b^2 = 7{,}3^2 - 4{,}8^2$ passt, und beschrifte die Seiten.}{\rechenplatz[halb]{4}}
\end{aufgabe}

\uebersichtskasten{Kathete: Ist die Hypotenuse bekannt, wird das Quadrat der anderen Kathete abgezogen – Gleichung umstellen, dann die Wurzel ziehen: a² = c² − b², a = √(c² − b²). \\ c = 13 cm, b = 5 cm: a² = 13² − 5² = 169 − 25 = 144 \quad a = √144 = 12 cm \\ Merkregel: Hypotenuse gesucht → plus. Kathete gesucht → minus. Die Kathete muss kürzer sein als die Hypotenuse, sonst stimmt etwas nicht. \\ Umkehrung: Gilt für die drei Seiten eines Dreiecks a² + b² = c² mit c als längster Seite, dann ist das Dreieck rechtwinklig; gilt es nicht, hat es keinen rechten Winkel. \\ 9 cm, 12 cm, 15 cm: 81 + 144 = 225 = 15² → rechtwinklig \quad 9 cm, 12 cm, 16 cm: 81 + 144 = 225 ≠ 256 → nicht rechtwinklig}
